\documentclass[../main.tex]{subfiles}

\begin{document}

\chapter*{Further Reading}
\addcontentsline{toc}{chapter}{Further Reading}
\label{app:further-reading}

This appendix provides resources for readers who wish to explore optimization for machine learning in greater depth. We organize recommendations by topic and difficulty level.

%% ============================================================
\section{Foundation: Mathematical Prerequisites}

Before diving deeper into optimization, ensure you have a solid foundation in the mathematical concepts from \emph{Book 1: The Math That Powers AI}:
\begin{itemize}
    \item Linear algebra: vectors, matrices, eigendecomposition, SVD
    \item Multivariable calculus: gradients, Hessians, Taylor series
    \item Probability and statistics: expectations, variance, random variables
\end{itemize}

%% ============================================================
\section{Textbooks on Convex Optimization}

\subsection{Essential Reading}

\textbf{Boyd, S. and Vandenberghe, L. (2004). \emph{Convex Optimization}. Cambridge University Press.}

The definitive textbook on convex optimization. Comprehensive coverage of convex sets, convex functions, duality theory, and algorithms. The book is freely available online at \texttt{stanford.edu/\textasciitilde boyd/cvxbook/}. Essential chapters for ML:
\begin{itemize}
    \item Chapter 2: Convex sets
    \item Chapter 3: Convex functions
    \item Chapter 5: Duality
    \item Chapter 9: Unconstrained minimization
    \item Chapter 10: Equality constrained minimization
\end{itemize}

\textbf{Nesterov, Y. (2018). \emph{Lectures on Convex Optimization} (2nd ed.). Springer.}

A more theoretical treatment by one of the field's pioneers. Covers optimal algorithms and complexity theory. Includes:
\begin{itemize}
    \item Accelerated gradient methods
    \item Optimal complexity bounds
    \item Smoothing techniques
    \item Interior point methods
\end{itemize}

\subsection{Additional Resources}

\textbf{Bertsekas, D. P. (2015). \emph{Convex Optimization Algorithms}. Athena Scientific.}

Focuses on algorithms with detailed convergence analysis. Good treatment of:
\begin{itemize}
    \item Proximal algorithms
    \item Alternating direction methods (ADMM)
    \item Incremental methods
\end{itemize}

%% ============================================================
\section{Numerical Optimization}

\textbf{Nocedal, J. and Wright, S. J. (2006). \emph{Numerical Optimization} (2nd ed.). Springer.}

The standard reference for numerical optimization algorithms. Comprehensive coverage of:
\begin{itemize}
    \item Line search and trust region methods
    \item Conjugate gradient methods
    \item Quasi-Newton methods (BFGS, L-BFGS)
    \item Constrained optimization (SQP, penalty methods)
    \item Large-scale optimization
\end{itemize}
Particularly useful for understanding classical optimization from which modern ML methods derive.

%% ============================================================
\section{Deep Learning Optimization}

\subsection{Textbooks}

\textbf{Goodfellow, I., Bengio, Y., and Courville, A. (2016). \emph{Deep Learning}. MIT Press.}

Chapter 8 (``Optimization for Training Deep Models'') provides excellent coverage of:
\begin{itemize}
    \item Challenges in neural network optimization
    \item Basic algorithms: SGD, momentum
    \item Adaptive learning rate methods
    \item Parameter initialization strategies
    \item Batch normalization
\end{itemize}
Available free online at \texttt{deeplearningbook.org}.

\textbf{Zhang, A., Lipton, Z. C., Li, M., and Smola, A. J. (2023). \emph{Dive into Deep Learning}. Cambridge University Press.}

Interactive textbook with runnable code. Available at \texttt{d2l.ai}. Relevant chapters:
\begin{itemize}
    \item Chapter 12: Optimization Algorithms
    \item Includes implementations in PyTorch, TensorFlow, and JAX
\end{itemize}

\subsection{Survey Papers}

\textbf{Ruder, S. (2016). ``An Overview of Gradient Descent Optimization Algorithms.'' arXiv:1609.04747.}

Excellent survey of modern optimization algorithms:
\begin{itemize}
    \item SGD variants and momentum
    \item AdaGrad, RMSprop, Adam
    \item Second-order methods
    \item Practical recommendations
\end{itemize}
Highly accessible; a great starting point for understanding different optimizers.

\textbf{Bottou, L., Curtis, F. E., and Nocedal, J. (2018). ``Optimization Methods for Large-Scale Machine Learning.'' \emph{SIAM Review}, 60(2), 223--311.}

Comprehensive survey bridging classical optimization and machine learning:
\begin{itemize}
    \item Stochastic gradient methods
    \item Noise reduction methods (mini-batching, variance reduction)
    \item Second-order methods for ML
    \item Parallel and distributed optimization
\end{itemize}

%% ============================================================
\section{Specialized Topics}

\subsection{Variance Reduction}

\textbf{Johnson, R. and Zhang, T. (2013). ``Accelerating Stochastic Gradient Descent using Predictive Variance Reduction.'' \emph{NeurIPS}.}

Introduces SVRG, achieving linear convergence for strongly convex problems.

\textbf{Defazio, A., Bach, F., and Lacoste-Julien, S. (2014). ``SAGA: A Fast Incremental Gradient Method with Support for Non-Strongly Convex Composite Objectives.'' \emph{NeurIPS}.}

Another variance reduction method with practical advantages.

\subsection{Adaptive Learning Rates}

\textbf{Kingma, D. P. and Ba, J. (2015). ``Adam: A Method for Stochastic Optimization.'' \emph{ICLR}.}

The original Adam paper. Essential reading for understanding:
\begin{itemize}
    \item Adaptive moment estimation
    \item Bias correction
    \item Default hyperparameter choices
\end{itemize}

\textbf{Loshchilov, I. and Hutter, F. (2019). ``Decoupled Weight Decay Regularization.'' \emph{ICLR}.}

Introduces AdamW, showing that weight decay and L2 regularization differ for adaptive optimizers.

\textbf{Reddi, S. J., Kale, S., and Kumar, S. (2018). ``On the Convergence of Adam and Beyond.'' \emph{ICLR}.}

Identifies convergence issues with Adam and proposes AMSGrad fix.

\subsection{Learning Rate Schedules}

\textbf{Smith, L. N. (2017). ``Cyclical Learning Rates for Training Neural Networks.'' \emph{WACV}.}

Introduces cyclical learning rates and learning rate range test.

\textbf{Loshchilov, I. and Hutter, F. (2017). ``SGDR: Stochastic Gradient Descent with Warm Restarts.'' \emph{ICLR}.}

Cosine annealing with warm restarts.

\textbf{Gotmare, A., Keskar, N. S., Xiong, C., and Socher, R. (2019). ``A Closer Look at Deep Learning Heuristics: Learning Rate Restarts, Warmup and Distillation.'' \emph{ICLR}.}

Analysis of practical training heuristics.

\subsection{Non-Convex Optimization}

\textbf{Ge, R., Huang, F., Jin, C., and Yuan, Y. (2015). ``Escaping From Saddle Points---Online Stochastic Gradient for Tensor Decomposition.'' \emph{COLT}.}

Shows how noise helps escape saddle points.

\textbf{Jin, C., Ge, R., Netrapalli, P., Kakade, S. M., and Jordan, M. I. (2017). ``How to Escape Saddle Points Efficiently.'' \emph{ICML}.}

Perturbed gradient descent for finding local minima.

%% ============================================================
\section{Online Courses}

\textbf{Stanford EE364a: Convex Optimization I}

Stephen Boyd's course based on his textbook. Lecture videos available online:
\begin{itemize}
    \item Theoretical foundations of convex optimization
    \item Duality theory
    \item Applications in signal processing and ML
\end{itemize}

\textbf{Stanford EE364b: Convex Optimization II}

Advanced topics including distributed optimization and applications.

\textbf{Berkeley CS 182: Deep Learning}

Covers practical aspects of training neural networks, including optimization.

\textbf{MIT 6.881: Optimization for Machine Learning}

Graduate-level course on optimization theory with ML focus.

%% ============================================================
\section{Software and Implementation}

\subsection{Optimization Libraries}

\begin{itemize}
    \item \textbf{PyTorch Optim}: \texttt{torch.optim} module with SGD, Adam, etc.
    \item \textbf{TensorFlow/Keras Optimizers}: \texttt{tf.keras.optimizers}
    \item \textbf{JAX Optax}: Composable gradient transformations
    \item \textbf{CVXPY}: Python library for convex optimization
    \item \textbf{SciPy Optimize}: Classical optimization algorithms
\end{itemize}

\subsection{Tutorials and Notebooks}

\begin{itemize}
    \item PyTorch tutorials on optimization: \texttt{pytorch.org/tutorials}
    \item JAX optimization examples: \texttt{github.com/google/jax}
    \item FastAI course materials: \texttt{course.fast.ai}
\end{itemize}

%% ============================================================
\section{Looking Ahead: Books 3--5}

This book provides the optimization foundation for the advanced topics in our series:

\textbf{Book 3: The Generative Revolution}
\begin{itemize}
    \item Optimization for recurrent networks (gradient clipping, truncated BPTT)
    \item Training transformers (warmup schedules, attention-specific techniques)
    \item Fine-tuning strategies for generative models
\end{itemize}

\textbf{Book 4: Inside the Black Box}
\begin{itemize}
    \item Variational inference as optimization
    \item Interpretability methods and their optimization
    \item Understanding loss landscape geometry
\end{itemize}

\textbf{Book 5: The Age of AI Agents}
\begin{itemize}
    \item Distributed optimization for multi-agent systems
    \item Meta-learning and learning to optimize
    \item Reinforcement learning optimization techniques
\end{itemize}

%% ============================================================
\section{Research Frontiers}

Active research areas in optimization for ML:

\begin{description}
\item[Sharpness-Aware Minimization (SAM)] Optimizing for flat minima to improve generalization.

\item[Large Batch Training] Techniques like LARS and LAMB for scaling to very large batch sizes.

\item[Learning Rate-Free Methods] Algorithms that automatically adapt without learning rate tuning (D-Adaptation, Prodigy).

\item[Optimizer Learning] Using neural networks to learn optimization algorithms.

\item[Continual Learning] Optimization without catastrophic forgetting.

\item[Federated Optimization] Distributed learning with communication constraints.

\item[Zeroth-Order Optimization] When gradients are unavailable (black-box optimization).
\end{description}

\bigskip

\noindent
\textbf{Note:} Web links are accurate as of 2024. For the most current versions and additional resources, search for the titles or authors mentioned.

\end{document}
